\documentclass[11pt]{article}
\usepackage[utf8]{inputenc}
\usepackage[T1]{fontenc}
\usepackage{amsmath,amssymb,mathtools,bm}
\usepackage{graphicx,booktabs,xcolor,array,longtable,float}
\usepackage[a4paper,margin=2.4cm]{geometry}
\usepackage[expansion=false]{microtype}
\usepackage[hidelinks]{hyperref}
\usepackage{caption}
\captionsetup{font=small,labelfont=bf}
\renewcommand{\figurename}{Figura}\renewcommand{\tablename}{Tabla}
\renewcommand{\abstractname}{Resumen}\renewcommand{\contentsname}{Índice}
\newcommand{\kPa}{\,\mathrm{kPa}}
\newcommand{\sbar}{\bar\sigma_{\theta\theta}}
\newcommand{\pend}{\textcolor{red!70!black}{\textbf{[pendiente]}}}
\setlength{\parindent}{0pt}\setlength{\parskip}{5pt}
\sloppy
\title{\textbf{Evaluación biomecánica transmural de la arteria carótida común:\\ formulación constitutiva, integración de pared gruesa,\\ modelo perivascular in vivo y comparación ex vivo}\\[4pt]\large Versión revisada y corregida}
\author{Simulación biomecánica y mecánica de medios continuos}
\date{Revisión del 6 de octubre de 2026 (sobre la versión del 30 de septiembre de 2026)}
\begin{document}
\maketitle
\begin{abstract}
Se presenta la formulación de mecánica de medios continuos para estimar la presión intraluminal $P_i$ en arterias carótidas comunes humanas, integrando los datos in vivo de Masson et al.\ (2011) con los ensayos ex vivo del complejo media--íntima de Sommer \& Holzapfel (2012). Respecto de la versión anterior, esta revisión: (i) define con precisión la tensión representativa $\sbar$ como el promedio ponderado de $(\sigma_{\theta\theta}-\sigma_{rr})$ sobre el espesor, con lo que la relación $P_{\rm pared}=\sbar\ln(1+h/r_i)$ resulta una identidad exacta; (ii) distingue el estiramiento $\lambda_\theta$ de Masson (con tensiones residuales) del estiramiento reducido $\lambda=r_i/R_i$ usado en las expresiones explícitas, y cuantifica la inconsistencia entre ambos en la Tabla~\ref{tab:geo}; (iii) corrige la bibliografía; (iv) reemplaza afirmaciones no sustentadas sobre la presión perivascular $P_a$ por un análisis de sensibilidad basado en los parámetros publicados por Masson et al.; y (v) unifica los intervalos de error de Laplace y de la aproximante de Padé $[1/1]$. Los resultados analíticos (sesgo de Laplace, Padé $[1/1]$, ley explícita $P_i(\lambda)$) se conservan y se verificaron numéricamente. Los valores que no pudieron contrastarse con las fuentes se marcan como \pend\ en la Sección~\ref{sec:pend}.
\end{abstract}
\tableofcontents

%--------------------------------------------------------------------
\section{Registro de correcciones respecto de la versión anterior}\label{sec:reg}
\begin{table}[H]\centering\small
\caption{Cambios introducidos. ``Fuente'' indica el documento contra el que se verificó.}
\begin{tabular}{p{3.0cm}p{4.6cm}p{4.9cm}p{2.0cm}}\toprule
Elemento & Versión anterior & Versión revisada & Fuente\\\midrule
Referencia Masson & ``A 3D constitutive model\ldots'', 10(4):511--522 & ``Carotid artery mechanical properties and stresses quantified using in vivo data from normotensive and hypertensive humans'', 10:867--882 & Archivo del artículo\\
Referencia Urquiza & ``(2012)'' & Tesis defendida en febrero de 2009; título ``Hemodinámica computacional del sistema arterial humano'' & Portada de la tesis\\
Referencia Sommer & título abreviado & título completo (``intact and layer-dissected'') & Cita en Masson (2010)\\
Definición de $\sbar$ & ``tensión de Cauchy circunferencial'' & promedio ponderado de $\sigma_{\theta\theta}-\sigma_{rr}$; la Ec.~(2) anterior pasa a ser identidad & Masson, Ec.~(11); Holzapfel et al.\ 2000, Ecs.~(27) y (31)\\
Estiramiento $\lambda_\theta$ & $r_i/R_i$ para todo el documento & Se distingue $\lambda_\theta(r)=\pi r/(\Theta_0R)$ (Masson) de $\lambda=r_i/R_i$ (modelo reducido) & Masson, Ecs.~(1)--(3)\\
Presión perivascular & $P_a=5.6$--$6.0\kPa$ constante, ``40\,\% de la presión'' & Modelo $P_a(t)=a\,e^{b\,r_m/r_m(t_d)}$; sensibilidad; valor adoptado marcado \pend & Masson, Ec.~(12), Tabla~4\\
Tensiones $\sbar$ de cohorte & presentadas como lecturas de Masson & se indica que no se pudieron verificar y que el texto de Masson da $61.9\pm18.4$ (HT) y $51.9\pm18.0\kPa$ (NT) & Masson, Sec.~3.2\\
Tabla~2 anterior (filas Sommer) & ``tensiones de Cauchy de Sommer'' & $\sbar$ obtenido de $P_i$ por inversión de la Ec.~(2): no es un dato experimental & Cálculo propio\\
Intervalos de error & $+6.7/+9.2$, $+6/+9$, $+4.92/+9.70\,\%$; Padé $<0.26$ y $<0.28\,\%$ & Un solo criterio: función de $x=h/r_i$ (Tabla~\ref{tab:pade}); en $x\in[0.10,0.20]$ Padé $\le0.28\,\%$ & Cálculo propio\\
Deformación axial & $\lambda_z\approx1$ & Se documenta que Masson estima $\lambda_z\approx1.10$ y se cuantifica el efecto (Tabla~\ref{tab:lz}) & Masson, Tabla~4\\
Ajuste ex vivo & ``tejido ex vivo'' en general & Donante VI~r, un solo espécimen, un $\lambda_z$ no especificado & Sommer 2012, Tabla~1, Fig.~6\\\bottomrule
\end{tabular}
\end{table}

%--------------------------------------------------------------------
\section{Fuentes primarias y figuras originales}
Las figuras originales (Fig.~5a y 7a de Masson et al.\ 2011; Fig.~6a de Sommer \& Holzapfel 2012) no se reproducen aquí porque la revisión se realizó sobre el texto de los documentos y no sobre sus imágenes; deben reinsertarse desde los artículos originales, con la autorización editorial correspondiente. Lo que sí se verificó en el texto de Masson es que ambas figuras expresan la tensión de Cauchy circunferencial \emph{en función del estiramiento en el radio medial $r_m$} (pie de las Figs.~5 y~7), que las curvas medias se calcularon cohorte por cohorte, y que las poblaciones son: 16 normotensos (NT, 21--64 años; 7 jóvenes de 21--42 y 9 mayores de 53--64) y 25 hipertensos tratados (HT, 44--69 años), con un subgrupo de 11 NT pareados por edad.

%--------------------------------------------------------------------
\section{Desarrollo físico, cinemático y constitutivo}
\subsection{Medición no invasiva}
Masson et al.\ miden con ecotracking (Art.Lab, Esaote; 30\,Hz, 6\,s) el diámetro medial y el espesor íntima--media (IMT), y con tonometría de aplanamiento (SphygmoCor; 128\,Hz, 9\,s) la presión en la misma carótida. La geometría y la presión se registran en instantes sucesivos ($5\pm2$ min de diferencia) y se alinean luego sobre un ciclo medio, tomando como referencia la fase de distensión sistólica; la geometría interna se deduce del diámetro medial y del IMT.

\subsection{Cinemática}\label{sec:cin}
\textbf{Modelo de Masson.} La pared se trata como un cilindro de pared gruesa con tensiones residuales, a partir de un sector abierto de ángulo $\Theta_0$. Con $J=\det\bm F=\lambda_r\lambda_\theta\lambda_z=1$, la deformación es
\begin{equation}
\lambda_\theta(r)=\frac{\pi\,r}{\Theta_0\,R},\qquad
R=\Bigl[R_m^2-\frac{\pi}{\Theta_0\lambda_z}\bigl(r_m^2-r^2\bigr)\Bigr]^{1/2},
\label{eq:masson}
\end{equation}
donde $R_m$ es el radio medial de referencia, $\Theta_0\approx130^\circ$ (Tabla~4 de Masson) y $\lambda_z$ el estiramiento axial, que en la Tabla~4 vale $\approx1.10$. Los estiramientos graficados en las Figs.~5 y 7 son los de $r=r_m$.

\textbf{Modelo reducido (sin tensiones residuales).} Para obtener expresiones explícitas se supone un tubo cargado a partir de la configuración sin carga, con $\lambda_z=1$ y estiramiento interno $\lambda=r_i/R_i$. La conservación del área de la sección de pared da
\begin{equation}
(r_i+h)^2-r_i^2=(R_i+H_0)^2-R_i^2\;\Longrightarrow\;
u(\lambda)\equiv1+\frac{h}{r_i}(\lambda)=\sqrt{1+\frac{(1+h_0)^2-1}{\lambda^{2}}},\quad h_0=\frac{H_0}{R_i}.
\label{eq:u}
\end{equation}
Si $\lambda_z\neq1$ la conservación del volumen sustituye $\lambda^2$ por $\lambda^2\lambda_z$ en \eqref{eq:u}.

\textbf{Observación.} $\lambda$ de \eqref{eq:u} y $\lambda_\theta(r_m)$ de \eqref{eq:masson} no son la misma variable. La Tabla~\ref{tab:geo} calcula, a partir de $r_i$, $h$ y $h_0$ de cada cohorte, los estiramientos interno $\lambda_i$ y medial $\lambda_m$ del modelo reducido: el $\lambda$ tabulado coincide con $\lambda_i$ para jóvenes, con $\lambda_m$ para mayores y queda entre ambos para HT. Esta diferencia es de hasta $0.03$--$0.04$ en $\lambda$, que, con el crecimiento exponencial de $\sbar(\lambda)$ en las tablas por cohorte (exponentes del orden de $8$--$10$), equivale a variaciones de la tensión del orden de $25$--$40\,\%$.

\subsection{Equilibrio radial y definición exacta de $\sbar$}
El equilibrio $\partial\sigma_{rr}/\partial r+(\sigma_{rr}-\sigma_{\theta\theta})/r=0$, con $\sigma_{rr}(r_i)=-P_i$ y $\sigma_{rr}(r_m)=-P_a$, integrado en el espesor da (Masson, Ec.~(11); Holzapfel et al.\ 2000, Ecs.~(27) y (31))
\begin{equation}
P_i=P_a+\int_{r_i}^{r_m}\frac{\sigma_{\theta\theta}(r)-\sigma_{rr}(r)}{r}\,dr .
\label{eq:eq}
\end{equation}
Se define la \emph{tensión circunferencial representativa}
\begin{equation}
\sbar\equiv\frac{1}{\ln(r_m/r_i)}\int_{r_i}^{r_m}\frac{\sigma_{\theta\theta}-\sigma_{rr}}{r}\,dr,
\label{eq:sbar}
\end{equation}
con lo que, \emph{sin aproximación adicional},
\begin{equation}
P_i=P_{\rm pared}+P_a,\qquad P_{\rm pared}=\sbar\,\ln\!\Bigl(1+\frac{h}{r_i}\Bigr),\qquad r_m=r_i+h .
\label{eq:Pi}
\end{equation}
La versión anterior llamaba a $\sbar$ ``tensión de Cauchy circunferencial'' sin precisar que es este promedio logarítmico de $\sigma_{\theta\theta}-\sigma_{rr}$; la identidad \eqref{eq:Pi} solo se sostiene con esa definición. $\sbar$ no coincide, en general, con $\sigma_{\theta\theta}$ evaluada en $r_m$, que es lo que grafican Masson et al.\ en las Figs.~5 y~7.

\subsection{Presión perivascular}
Masson et al.\ combinan el efecto del tejido perivascular y de la adventicia en un único término que varía con el ciclo cardíaco (su Ec.~(12), citando a Humphrey y Na, 2002):
\begin{equation}
P_a(t)=a\,\exp\!\Bigl(b\,\frac{r_m(t)}{r_m(t_d)}\Bigr),
\label{eq:Pa}
\end{equation}
donde $t_d$ es el instante diastólico. Para los 25 pacientes HT (Tabla~4), la media es $a=0.078\pm0.036\kPa$ y $b=3.56\pm0.06$. Evaluado con la Ec.~\eqref{eq:Pa} \emph{leída literalmente} (el texto extraído del PDF puede degradar la fórmula; conviene contrastarla con el original) y con los $a$, $b$ individuales:
\begin{center}\small
\begin{tabular}{lcccc}\toprule
$r_m/r_m(t_d)$ & Media [kPa] & DE [kPa] & Mínimo [kPa] & Máximo [kPa]\\\midrule
1.000 & 2.76 & 1.30 & 0.14 & 4.91\\
1.034 & 3.11 & 1.47 & 0.16 & 5.55\\
1.060 & 3.41 & 1.61 & 0.17 & 6.09\\\bottomrule
\end{tabular}
\end{center}
El valor constante $P_a=5.6$--$6.0\kPa$ de la versión anterior queda así en el extremo superior del rango individual y no es representativo de la media de la cohorte HT; su origen no está documentado \pend. Además, \eqref{eq:Pa} varía en el ciclo, mientras que las tablas por cohorte (Sec.~\ref{sec:tablas}) lo tratan como constante. Se mantiene el valor anterior en las tablas para conservar la continuidad con el resto del análisis, pero las conclusiones sobre $P_a$ se reformulan como dependientes de él (Sec.~\ref{sec:sens}).

\subsection{Sesgo de la ley de Laplace (pared delgada)}
En la tesis de Urquiza (Cap.~3, Ec.~3.13) la pared delgada se escribe $(p-p_{\rm ext})R=\sigma h$, es decir, $P_{\rm pared}=\sigma\,h/R$; Holzapfel et al.\ (2000) observan que esta aproximación de membrana es incompatible con tensiones residuales. Evaluada en $R=r_i$ y con $x=h/r_i$,
\begin{equation}
\frac{P_{\rm Laplace}}{P_{\rm exacta}}=\frac{x}{\ln(1+x)}=1+\frac{x}{2}-\frac{x^{2}}{12}+\frac{x^{3}}{24}-\cdots
\;\Longrightarrow\;\text{sesgo}\approx\frac{x}{2}\times100\,\%\ (\text{orden principal}).
\label{eq:laplace}
\end{equation}
Con $x=0.138$--$0.189$ (cohortes de la Tabla~\ref{tab:geo}) el sesgo exacto es $+6.7$ a $+9.1\,\%$, compatible con el intervalo ``$+6.7\,\%$ a $+9.2\,\%$'' de la versión anterior. Sobre todo el rango fisiológico $x\in[0.10,0.20]$ es $+4.92$ a $+9.70\,\%$ (Tabla~\ref{tab:pade}). Ambos intervalos son correctos: responden a dominios distintos, y no deben mezclarse en el texto.

\subsection{Formulación de Urquiza evaluada en el radio medio}
Con $R=r_m^{\rm mid}=r_i+h/2$,
\begin{equation}
P_{{\rm Urq},r_{\rm mid}}=\sbar\,\frac{h/r_i}{1+\tfrac12\,h/r_i}=\sbar\,\frac{2x}{2+x}=\sbar\,\frac{2(u-1)}{u+1},
\label{eq:urq}
\end{equation}
que es la aproximante de Padé $[1/1]$ de $\ln u$ en torno a $u=1$. El desarrollo da $2x/(2+x)=x-\tfrac{x^2}{2}+\tfrac{x^3}{4}-\cdots$ frente a $\ln(1+x)=x-\tfrac{x^2}{2}+\tfrac{x^3}{3}-\cdots$, por lo que el error relativo es $\approx-x^2/12$, siempre negativo. Su valor numérico (Tabla~\ref{tab:pade}) llega a $-0.276\,\%$ en $x=0.20$; en $x\in[0.10,0.20]$ el error de Padé no supera $0.28\,\%$ en valor absoluto. (Las menciones previas a ``$<0.26\,\%$'' valen solo hasta $x\approx0.19$.)

\begin{table}[H]\centering\small
\caption{Factor logarítmico exacto $\ln u$, aproximante de Padé $[1/1]$ y ley de Laplace ($u-1$), con $u=1+h/r_i$. Todos los valores se recalcularon; coinciden con la Tabla~1 de la versión anterior.}\label{tab:pade}
\begin{tabular}{ccccccc}\toprule
$x=h/r_i$ & $u$ & $\ln u$ & Padé & Err.\ Padé & $u-1$ & Err.\ Laplace\\\midrule
0.00 & 1.00 & 0.00000 & 0.00000 & +0.000\,\% & 0.00000 & +0.00\,\% \\
0.05 & 1.05 & 0.04879 & 0.04878 & -0.020\,\% & 0.05000 & +2.48\,\% \\
0.10 & 1.10 & 0.09531 & 0.09524 & -0.076\,\% & 0.10000 & +4.92\,\% \\
0.12 & 1.12 & 0.11333 & 0.11321 & -0.107\,\% & 0.12000 & +5.89\,\% \\
0.15 & 1.15 & 0.13976 & 0.13953 & -0.162\,\% & 0.15000 & +7.33\,\% \\
0.18 & 1.18 & 0.16551 & 0.16514 & -0.228\,\% & 0.18000 & +8.75\,\% \\
0.20 & 1.20 & 0.18232 & 0.18182 & -0.276\,\% & 0.20000 & +9.70\,\% \\
0.25 & 1.25 & 0.22314 & 0.22222 & -0.413\,\% & 0.25000 & +12.04\,\% \\

\bottomrule\end{tabular}
\end{table}
\begin{figure}[H]\centering\includegraphics[width=.62\linewidth]{fig/f_pade.pdf}
\caption{Arriba: factor logarítmico exacto, Padé $[1/1]$ y Laplace. Abajo: error relativo respecto de $\ln(1+x)$. La banda dorada es el rango fisiológico $x\in[0.10,0.20]$.}\end{figure}

\subsection{Leyes explícitas $P_i(\lambda)$}
Sustituyendo \eqref{eq:u} en \eqref{eq:Pi} y usando $\ln u=\tfrac12\ln u^2$:
\begin{align}
P_{i,\rm exacta}(\lambda)&=\tfrac12\,\sbar(\lambda)\,\ln\!\Bigl[1+\tfrac{(1+h_0)^2-1}{\lambda^{2}}\Bigr]+P_a,\label{eq:Pexacta}\\
P_{i,\rm Urq}(\lambda)&=\sbar(\lambda)\,\frac{2(u-1)}{u+1}+P_a,\qquad u=\sqrt{1+\tfrac{(1+h_0)^2-1}{\lambda^{2}}}.\label{eq:Purq}
\end{align}
Estas expresiones se verificaron numéricamente (Sec.~\ref{sec:tablas}).

\subsection{Relación entre las tensiones de Masson y las tensiones de las Tablas~3--6}\label{sec:rel}
Las trayectorias de Masson et al.\ no son ensayos de tracción: son estimaciones indirectas obtenidas ajustando por mínimos cuadrados 14 parámetros (residuales, pasivos, activos y perivasculares) a la presión medida (Masson, Sec.~2.2.4; función \texttt{lsqnonlin}). Las tensiones que muestran las Figs.~5 y~7 son $\sigma_{\theta\theta}(r_m)$ a lo largo del ciclo. Los valores $\sbar$ de las tablas por cohorte son, en cambio, una tensión representativa en el sentido de \eqref{eq:sbar}, y su procedencia numérica no pudo reconstruirse (Sec.~\ref{sec:pend}). Despejando \eqref{eq:Pi},
\begin{equation}
\sbar(\lambda)=\frac{P_i(\lambda)-P_a}{\ln\bigl(1+h(\lambda)/r_i(\lambda)\bigr)} .
\label{eq:inv}
\end{equation}
\textbf{Aclaración física.} Para $\sbar=60\kPa$ en la cohorte joven, $P_{\rm pared}=7.8\kPa$ ($58.6$ mmHg), y con $P_a=5.63\kPa$ resulta $P_i=13.44\kPa$ ($100.8$ mmHg). La partición entre pared y entorno perivascular depende de $P_a$ (Sec.~\ref{sec:sens}).

%--------------------------------------------------------------------
\section{Ajuste analítico ex vivo}
Para el complejo media--íntima de la carótida común del donante VI~r (Sommer \& Holzapfel 2012, Fig.~6a) se ajustó
\begin{equation}
P_i(\lambda_\theta)=0.1427\,\bigl(e^{\,11.9103(\lambda_\theta-1)}-1\bigr)\ \text{[kPa]},
\label{eq:exvivo}
\end{equation}
que vale $1.40$, $4.94$ y $16.59\kPa$ en $\lambda_\theta=1.2$, $1.3$ y $1.4$ (la versión anterior daba $1.43$, $5.12$ y $17.09\kPa$, un $2$--$3\,\%$ más; esas filas deben considerarse lecturas de la figura y no del ajuste). Alcanza $P_0=13.32\kPa$ (100 mmHg) en $\lambda_\theta=1.382$. Conviene acotar el alcance:
\begin{itemize}
\item Es \emph{un solo espécimen}. En la Tabla~1 de Sommer \& Holzapfel, el media--íntima de VI~r tiene $k_2=2.8$, frente a una media de $16.5$ (DE $9.3$) entre los 4 media--íntima de carótida común; el resto de sus parámetros es típico ($\mu=118.9\kPa$ frente a $122.3\pm12.2$).
\item Los ensayos de media--íntima se hicieron a presiones de 0--20\,kPa y estiramientos axiales de 0 a 15\,\%; el $\lambda_z$ del ajuste no está indicado \pend. Sommer y Holzapfel señalan que la mecánica de los tubos media--íntima es casi independiente del estiramiento axial aplicado, lo que atenúa, sin eliminar, esta ambigüedad.
\item $P_i(1)=0$ está impuesto por la forma \eqref{eq:exvivo}; el término $-1$ obliga a la curva a pasar por el origen y no es un resultado del ensayo.
\item Sommer y Holzapfel usan el marco de Holzapfel et al.\ (2000), donde $\lambda_{\theta i}$ es el estiramiento de la superficie interna e incluye las deformaciones residuales (parámetro $k$ del ángulo de apertura); su equivalencia con el $\lambda$ de Masson o con el modelo reducido \eqref{eq:u} no está establecida \pend.
\end{itemize}
Las filas ``Sommer'' de la Tabla~2 de la versión anterior calculaban $\sbar$ despejando de $P_i$ mediante \eqref{eq:inv} con $P_a=0$: no son tensiones medidas por Sommer \& Holzapfel. Las geometrías $r_i$, $h$ de esas filas tampoco figuran en el artículo y se omiten.

\textbf{Contexto adicional.} Valdez-Jasso et al.\ (2011) comparan in vivo y ex vivo en ovino y en carótidas humanas, y encuentran que los vasos son más rígidos ex vivo (para un aumento de 100\,mmHg, la sección de la carótida ovina se distiende $\approx2\,\%$ ex vivo y la de la carótida humana $\approx18\,\%$ in vivo; son especies distintas). Lo atribuyen posiblemente a la extracción del vaso de su matriz de soporte. Su modelo es de pared delgada (Laplace con $r/h$), con las limitaciones de \eqref{eq:laplace}.

%--------------------------------------------------------------------
\section{Cuantificación por cohorte}\label{sec:tablas}
\subsection{Tabla resumen: consistencia geométrica}
\begin{table}[H]\centering\small
\caption{Puntos fisiológicos de la versión anterior (Tabla~2 de esa versión) y su consistencia. $\lambda_i,\lambda_m$: estiramientos interno y medial del modelo reducido calculados con la geometría de la fila (conservación de área y $h_0$ de la cohorte). $x_{\rm geo}=h/r_i$ medido; $x_{\rm cin}$ de \eqref{eq:u}. Columnas de $P_{\rm pared}$: tabulado, con $x_{\rm geo}$ y con $x_{\rm cin}$ [kPa].}\label{tab:geo}
\setlength{\tabcolsep}{4pt}
\begin{tabular}{lccccccccc}\toprule
Cohorte & $\lambda$ tab. & $\lambda_i$ & $\lambda_m$ & $x_{\rm geo}$ & $x_{\rm cin}$ & $P_{\rm pared}$ tab. & con $x_{\rm geo}$ & con $x_{\rm cin}$\\\midrule
Jóvenes & 1.115 & 1.120 & 1.089 & 0.138 & 0.139 & 7.79 & 7.75 & 7.81 \\
Mayores & 1.092 & 1.133 & 1.093 & 0.164 & 0.176 & 11.15 & 10.96 & 11.69 \\
Hipertensos & 1.075 & 1.087 & 1.059 & 0.189 & 0.192 & 14.42 & 14.17 & 14.43 \\

\bottomrule\end{tabular}
\end{table}
Las diferencias entre las dos últimas columnas son pequeñas para jóvenes e HT y de $0.7\kPa$ ($6\,\%$) para mayores, cuyo $P_{\rm pared}$ de la versión anterior ($11.15\kPa$ en la Tabla~2) no coincide con ninguna de las dos estimaciones; la Tabla~4 anterior da $11.69\kPa$, que corresponde a $x_{\rm cin}$. Debe unificarse el criterio \pend.

\subsection{Tablas por cohorte ($\lambda\in[1.000,1.180]$)}
Con $P_0=13.32\kPa\equiv100$ mmHg. Los valores $\sbar$ son los de la versión anterior; todo lo demás se recalculó con \eqref{eq:u}, \eqref{eq:Pi}, \eqref{eq:Pexacta} y \eqref{eq:Purq}. El máximo de $|P_{i,\rm recalc}-P_{i,\rm anterior}|$ es $0.003$--$0.005\kPa$ en las cuatro cohortes, es decir, las tablas anteriores son internamente consistentes; el error de Urquiza es el de $P_i$ total (incluye $P_a$).
\def\tablehead{$\lambda$ & $\sbar$ [kPa] & $h/r_i$ & $P_{\rm pared}$ [kPa] & $P_i$ [kPa] & $P_i$ [mmHg] & $P_i/P_0$ & $P_{i,\rm Urq}$ [kPa] & Err.\\\midrule}
\begin{table}[H]\centering\footnotesize\setlength{\tabcolsep}{3pt}
\caption{Jóvenes normotensos (Mean Young, 21--42 años). $h_0=0.1703$, $P_a=5.63\kPa$ \pend. Punto medido: $\lambda=1.115$, $\sbar=60.0\kPa$.}
\begin{tabular}{ccccccccc}\toprule\tablehead1.000 & 0.00 & 0.170 & 0.00 & 5.63 & 42.2 & 0.423 & 5.63 & +0.00\,\% \\
1.020 & 5.45 & 0.164 & 0.83 & 6.46 & 48.4 & 0.485 & 6.46 & -0.02\,\% \\
1.040 & 13.00 & 0.158 & 1.91 & 7.54 & 56.6 & 0.566 & 7.54 & -0.05\,\% \\
1.060 & 22.65 & 0.153 & 3.22 & 8.85 & 66.4 & 0.664 & 8.85 & -0.06\,\% \\
1.080 & 34.39 & 0.148 & 4.73 & 10.36 & 77.7 & 0.778 & 10.36 & -0.07\,\% \\
1.100 & 48.24 & 0.143 & 6.43 & 12.06 & 90.5 & 0.905 & 12.05 & -0.08\,\% \\
1.115 & 60.00 & 0.139 & 7.81 & 13.44 & 100.8 & 1.009 & 13.43 & -0.08\,\% \\
1.120 & 64.18 & 0.138 & 8.29 & 13.92 & 104.4 & 1.045 & 13.91 & -0.08\,\% \\
1.140 & 82.23 & 0.133 & 10.29 & 15.92 & 119.4 & 1.195 & 15.91 & -0.08\,\% \\
1.160 & 102.37 & 0.129 & 12.42 & 18.05 & 135.4 & 1.355 & 18.04 & -0.08\,\% \\
1.180 & 124.61 & 0.125 & 14.67 & 20.30 & 152.2 & 1.524 & 20.28 & -0.08\,\% \\
\bottomrule\end{tabular}\end{table}
\begin{table}[H]\centering\footnotesize\setlength{\tabcolsep}{3pt}
\caption{Mayores normotensos (Mean Old, 53--64 años). $h_0=0.2072$, $P_a=5.61\kPa$ \pend. Punto medido: $\lambda=1.092$, $\sbar=72.0\kPa$.}
\begin{tabular}{ccccccccc}\toprule\tablehead1.000 & 0.00 & 0.207 & 0.00 & 5.61 & 42.1 & 0.421 & 5.61 & +0.00\,\% \\
1.020 & 7.32 & 0.200 & 1.33 & 6.94 & 52.1 & 0.521 & 6.94 & -0.05\,\% \\
1.040 & 19.26 & 0.193 & 3.40 & 9.01 & 67.6 & 0.676 & 9.00 & -0.10\,\% \\
1.060 & 35.84 & 0.186 & 6.12 & 11.73 & 88.0 & 0.881 & 11.71 & -0.13\,\% \\
1.080 & 57.05 & 0.180 & 9.44 & 15.05 & 112.9 & 1.130 & 15.02 & -0.14\,\% \\
1.092 & 72.00 & 0.176 & 11.69 & 17.30 & 129.7 & 1.299 & 17.27 & -0.15\,\% \\
1.100 & 82.89 & 0.174 & 13.29 & 18.90 & 141.7 & 1.419 & 18.87 & -0.15\,\% \\
1.120 & 113.36 & 0.168 & 17.62 & 23.23 & 174.2 & 1.744 & 23.19 & -0.15\,\% \\
1.140 & 148.47 & 0.163 & 22.38 & 27.99 & 210.0 & 2.102 & 27.95 & -0.15\,\% \\
1.160 & 188.20 & 0.158 & 27.53 & 33.14 & 248.6 & 2.488 & 33.09 & -0.15\,\% \\
1.180 & 232.57 & 0.153 & 33.03 & 38.64 & 289.8 & 2.901 & 38.58 & -0.14\,\% \\
\bottomrule\end{tabular}\end{table}
\begin{table}[H]\centering\footnotesize\setlength{\tabcolsep}{3pt}
\caption{Normotensos de control (Mean NT, 39--64 años). $h_0=0.1852$, $P_a=5.62\kPa$ \pend. Punto medido: $\lambda=1.095$, $\sbar=65.0\kPa$.}
\begin{tabular}{ccccccccc}\toprule\tablehead1.000 & 0.00 & 0.185 & 0.00 & 5.62 & 42.2 & 0.422 & 5.62 & +0.00\,\% \\
1.020 & 6.67 & 0.179 & 1.10 & 6.72 & 50.4 & 0.504 & 6.71 & -0.04\,\% \\
1.040 & 17.08 & 0.172 & 2.71 & 8.33 & 62.5 & 0.626 & 8.33 & -0.07\,\% \\
1.060 & 31.23 & 0.166 & 4.80 & 10.42 & 78.2 & 0.783 & 10.41 & -0.09\,\% \\
1.080 & 49.13 & 0.161 & 7.32 & 12.94 & 97.0 & 0.971 & 12.92 & -0.10\,\% \\
1.095 & 65.00 & 0.157 & 9.45 & 15.07 & 113.0 & 1.132 & 15.06 & -0.11\,\% \\
1.100 & 70.76 & 0.155 & 10.21 & 15.83 & 118.7 & 1.188 & 15.81 & -0.11\,\% \\
1.120 & 96.13 & 0.150 & 13.44 & 19.06 & 143.0 & 1.431 & 19.04 & -0.11\,\% \\
1.140 & 125.25 & 0.145 & 16.98 & 22.60 & 169.5 & 1.697 & 22.57 & -0.11\,\% \\
1.160 & 158.10 & 0.141 & 20.79 & 26.41 & 198.1 & 1.982 & 26.38 & -0.11\,\% \\
1.180 & 194.70 & 0.136 & 24.84 & 30.46 & 228.5 & 2.287 & 30.42 & -0.11\,\% \\
\bottomrule\end{tabular}\end{table}
\begin{table}[H]\centering\footnotesize\setlength{\tabcolsep}{3pt}
\caption{Hipertensos tratados (Mean HT, 44--69 años). $h_0=0.2197$, $P_a=5.98\kPa$ \pend. Punto medido: $\lambda=1.075$, $\sbar=82.0\kPa$.}
\begin{tabular}{ccccccccc}\toprule\tablehead1.000 & 0.00 & 0.220 & 0.00 & 5.98 & 44.9 & 0.449 & 5.98 & +0.00\,\% \\
1.020 & 10.23 & 0.212 & 1.97 & 7.95 & 59.6 & 0.597 & 7.94 & -0.08\,\% \\
1.040 & 28.92 & 0.205 & 5.38 & 11.36 & 85.2 & 0.853 & 11.35 & -0.14\,\% \\
1.060 & 56.08 & 0.198 & 10.11 & 16.09 & 120.7 & 1.208 & 16.06 & -0.17\,\% \\
1.075 & 82.00 & 0.192 & 14.43 & 20.41 & 153.1 & 1.533 & 20.38 & -0.18\,\% \\
1.080 & 91.70 & 0.191 & 16.02 & 22.00 & 165.0 & 1.651 & 21.96 & -0.18\,\% \\
1.100 & 135.78 & 0.184 & 22.99 & 28.97 & 217.3 & 2.175 & 28.92 & -0.19\,\% \\
1.120 & 188.32 & 0.178 & 30.92 & 36.90 & 276.8 & 2.771 & 36.83 & -0.19\,\% \\
1.140 & 249.32 & 0.173 & 39.72 & 45.70 & 342.8 & 3.431 & 45.62 & -0.18\,\% \\
1.160 & 318.79 & 0.167 & 49.29 & 55.27 & 414.6 & 4.150 & 55.18 & -0.18\,\% \\
1.180 & 396.72 & 0.162 & 59.56 & 65.54 & 491.6 & 4.921 & 65.43 & -0.17\,\% \\
\bottomrule\end{tabular}\end{table}

\subsection{Sensibilidad a $P_a$ y a la deformación axial}\label{sec:sens}
Si la presión total $P_i$ del punto fisiológico se mantiene fija (Tabla~2 anterior), la tensión representativa necesaria depende de $P_a$ mediante \eqref{eq:inv}:
\begin{table}[H]\centering\small
\caption{$\sbar$ [kPa] requerido para reproducir $P_i$ según el $P_a$ adoptado. $P_a$: valor de la versión anterior; $3.10$ y $2.75\kPa$: valores de Masson (media HT, $r_m/r_m(t_d)=1.034$ y $1.0$); $0$: sin soporte perivascular.}\label{tab:sens}
\begin{tabular}{lcccccc}\toprule
Cohorte & $P_i$ [kPa] & $\sbar$ tab. & $P_a$ anterior & $3.10$ & $2.75$ & $0$\\\midrule
Jóvenes & 13.42 & 60.0 & 59.9 & 79.3 & 82.0 & 103.1 \\
Mayores & 16.76 & 72.0 & 68.7 & 84.2 & 86.3 & 103.3 \\
Hipertensos & 20.40 & 82.0 & 81.9 & 98.3 & 100.3 & 115.9 \\
\bottomrule\end{tabular}
\end{table}
Con los valores de Masson, la tensión representativa necesaria sube entre $+20$ y $+32\,\%$ respecto de la tabla (y entre $+41$ y $+72\,\%$ si se anulara $P_a$); por tanto, el análisis de $\sbar$ por cohorte y la afirmación de que el entorno perivascular ``asume cerca del 40\,\% de la presión'' dependen de un parámetro sin respaldo documental. Con los valores de Masson, la fracción soportada por el entorno es $\approx15\,\%$ de $P_i$ en HT ($3.1/20.4$).

Por otro lado, Masson estima $\lambda_z\approx1.10$ para los 25 HT; con $\lambda_z=1.10$ en lugar de $1.0$ en \eqref{eq:u}, el espesor relativo y la presión de pared cambian como:
\begin{table}[H]\centering\small
\caption{Efecto de $\lambda_z=1.10$ (vs.\ $1.00$) en el punto fisiológico con $\sbar$ fijo.}\label{tab:lz}
\begin{tabular}{lccccc}\toprule
Cohorte & $x$ ($\lambda_z=1$) & $x$ ($\lambda_z=1.1$) & $P_{\rm pared}$ ($\lambda_z=1$) & $P_{\rm pared}$ ($\lambda_z=1.1$) & Variación\\\midrule
Jóvenes & 0.139 & 0.127 & 7.81 & 7.18 & -8.1\,\% \\
Mayores & 0.176 & 0.161 & 11.69 & 10.77 & -7.9\,\% \\
Hipertensos & 0.192 & 0.176 & 14.43 & 13.31 & -7.8\,\% \\
\bottomrule\end{tabular}
\end{table}
La reducción es de $\approx8\,\%$ en $P_{\rm pared}$, comparable al sesgo de Laplace que el documento intenta corregir; por eso la hipótesis $\lambda_z\approx1$ debe declararse. Los valores de $h_0$ por cohorte, además, se definieron respecto de una configuración de referencia que no se especifica \pend.

%--------------------------------------------------------------------
\section{Superposición y escala normalizada}
La Figura~\ref{fig:sup} superpone las curvas in vivo (calculadas con las tablas por cohorte) y la curva ex vivo \eqref{eq:exvivo}, en kPa y normalizadas por $P_0=13.32\kPa$.
\begin{figure}[H]\centering\includegraphics[width=\linewidth]{fig/f_super.pdf}
\caption{Izquierda: $P_i(\lambda_\theta)$. Derecha: $P_i/P_0$ (la línea discontinua es $P/P_0=1$). Las curvas in vivo usan $P_a$ de la versión anterior; el ex vivo es un solo espécimen y su $\lambda_\theta$ no está definido necesariamente igual (Sec.~4).}\label{fig:sup}
\end{figure}
Con los valores anteriores, la arteria in vivo alcanza $P/P_0=1$ (interpolando linealmente las tablas) en $\lambda\approx1.114$ (jóvenes), $1.083$ (NT), $1.070$ (mayores) y $1.048$ (HT), mientras que el ex vivo requiere $\lambda_\theta=1.382$. Una parte de esa diferencia es la del origen: $P_i(1)=P_a\approx5.6$--$6.0\kPa$ in vivo frente a $0$ ex vivo.

%--------------------------------------------------------------------
\section{Conclusiones revisadas}
\begin{enumerate}
\item \textbf{Formulación.} Con $\sbar$ definido por \eqref{eq:sbar}, $P_i=\sbar\ln(1+h/r_i)+P_a$ es exacta. El sesgo de Laplace es $x/\ln(1+x)-1\approx x/2$: $+4.9$ a $+9.7\,\%$ para $x\in[0.10,0.20]$, y $+6.7$ a $+9.1\,\%$ en las tres cohortes de la Tabla~\ref{tab:geo}.
\item \textbf{Urquiza.} Evaluar en $R=r_i+h/2$ equivale a la aproximante de Padé $[1/1]$; su error relativo $\approx-x^2/12$ es $\le0.28\,\%$ en $[0.10,0.20]$ (dato recalculado).
\item \textbf{Perivascular.} El análisis de Masson contradice un $P_a$ constante: $P_a(t)$ varía en el ciclo y su media en HT es $\approx2.8$--$3.1\kPa$ (individuos $0.1$--$5.6\kPa$). La proporción de $P_i$ que soporta el entorno no puede afirmarse sin fijar $P_a$ \pend.
\item \textbf{Normalizada.} El ex vivo requiere $\lambda_\theta=1.382$ frente a $1.05$--$1.11$ in vivo (según cohorte) para $P/P_0=1$; es un resultado para un espécimen y bajo definiciones de $\lambda$ no necesariamente equivalentes.
\item \textbf{Alcance.} Las tablas por cohorte son internamente consistentes, pero su ingreso ($\sbar$, $P_a$, $h_0$, $\lambda$) debe verificarse contra Masson antes de usarse como ``datos'' (p.\,ej.\ para entrenar modelos).
\end{enumerate}

%--------------------------------------------------------------------
\section{Elementos que no se pudieron verificar}\label{sec:pend}
\begin{enumerate}
\item Origen de $P_a=5.63$, $5.61$, $5.62$, $5.98\kPa$ (no figuran en la Tabla~4 de Masson; su media de $a$, $b$ da $\approx2.8$--$3.1\kPa$). La Tabla~2 de Masson (NT) no estaba disponible para comparar.
\item Origen de $\sbar=60$, $72$, $65$, $82\kPa$ y de las funciones $\sbar(\lambda)$ (lectura de figura, ajuste exponencial o cálculo): el texto de Masson da $61.9\pm18.4$ (HT) y $51.9\pm18.0\kPa$ (NT) en el estiramiento medio de la pared externa, que no coincide con $82$ y $65\kPa$ salvo que se trate de $\lambda$ distintos; $+14\,\%$ de mayores sobre jóvenes frente a $+20\,\%$ de las tablas.
\item Los $h_0=0.1703$, $0.2072$, $0.1852$, $0.2197$ y la configuración de referencia a la que se refieren.
\item El $\lambda_z$ del ajuste ex vivo \eqref{eq:exvivo} y el origen de las filas $\lambda=1.2$, $1.3$, $1.4$ de la Tabla~2 anterior.
\item Las figuras originales (Sec.~2) y los datos bibliográficos marcados con $^\dagger$.
\end{enumerate}
Además, el análisis de $P_a$ usa la Ec.~(12) de Masson tal como la entrega el texto extraído del PDF; hay que confirmarla contra el original.

%--------------------------------------------------------------------
\begin{thebibliography}{9}\small
\bibitem{Masson2011} Masson I., Beaussier H., Boutouyrie P., Laurent S., Humphrey J.D., Zidi M. (2011). Carotid artery mechanical properties and stresses quantified using in vivo data from normotensive and hypertensive humans. \emph{Biomechanics and Modeling in Mechanobiology}, 10:867--882. doi:10.1007/s10237-010-0279-6.
\bibitem{Sommer2012} Sommer G., Holzapfel G.A. (2012). 3D constitutive modeling of the biaxial mechanical response of intact and layer-dissected human carotid arteries. \emph{Journal of the Mechanical Behavior of Biomedical Materials}, 5:116--128.$^\dagger$
\bibitem{Sommer2010} Sommer G., Regitnig P., Költringer L., Holzapfel G.A. (2010). Biaxial mechanical properties of intact and layer-dissected human carotid arteries at physiological and supraphysiological loadings. \emph{Am.\ J.\ Physiol.\ Heart Circ.\ Physiol.}, 298:H898--H912.
\bibitem{Holz2000} Holzapfel G.A., Gasser T.C., Ogden R.W. (2000). A new constitutive framework for arterial wall mechanics and a comparative study of material models. \emph{Journal of Elasticity}, 61:1--48.
\bibitem{HolzOgd2010} Holzapfel G.A., Ogden R.W. (2010). Constitutive modelling of arteries. \emph{Proc.\ R.\ Soc.\ A}, 466:1551--1597.$^\dagger$
\bibitem{HolzGasser} Holzapfel G.A., Gasser T.C. Structural and numerical models for the (visco)elastic response of arterial walls with residual stresses. En: Holzapfel G.A., Ogden R.W. (eds.), \emph{Biomechanics of Soft Tissue in Cardiovascular Systems}, CISM 441, Springer.$^\dagger$
\bibitem{Valdez2011} Valdez-Jasso D., Bia D., Zócalo Y., Armentano R.L., Haider M.A., Olufsen M.S. (2011). Linear and nonlinear viscoelastic modeling of aorta and carotid pressure--area dynamics under in vivo and ex vivo conditions. \emph{Annals of Biomedical Engineering}, 39(5):1438--1456.
\bibitem{Urquiza2009} Urquiza S.A. (2009). \emph{Hemodinámica computacional del sistema arterial humano}. Tesis de Doctorado en Ciencias de la Ingeniería, Instituto Balseiro, Universidad Nacional de Cuyo--CNEA, San Carlos de Bariloche.
\bibitem{Caballero2020} Caballero D.E. (2020). \emph{Modelado constitutivo multiescala de matrices nanofibrosas electrohiladas}. Tesis doctoral, Universidad Nacional de Mar del Plata (director: S.A. Urquiza). Mencionada como contexto (compliancia de injertos vasculares); no se usa para ningún resultado de este documento.
\end{thebibliography}
$^\dagger$ Datos de edición no verificados contra el artículo (solo contra el índice de citas de otro documento o de memoria).
\end{document}
